\subsection{Two Lemmas on Bimodules}

Let A and B be rings. For any homomorphism f from B to A, we denote by \(A^{f}\) the (B,A)-bimodule with underlying right A-module \(A_{d}\) and external law for its left B-module structure given by \((b,a)\mapsto f(b)a\) .

Lemma 1. --- Let f and g be homomorphisms from B to A. The following conditions are equivalent:

\begin{enumerate}
\def\labelenumi{(\roman{enumi})}
\item
  The (B,A)-bimodules \(A^{f}\) and \(A^{g}\) are isomorphic.
\item
  There exists an inner automorphism (I, §8, No.~4, p.~102, Example 2) \(\theta\) of A such that \(g = \theta \circ f\) .
\end{enumerate}

The automorphisms of the right A-module \(A_{d}\) are the mappings \(x \mapsto ax\) , where a is an invertible element of A. Such an automorphism is a B-linear mapping from \(A^{f}\) to \(A^{g}\) if and only if we have

\[
g (b) a x = a f (b) x
\]

for every x in A and every b in B. This relation is equivalent to \(g(b) = af(b)a^{-1}\) for every b in B, that is, to \(g = \theta \circ f\) , where \(\theta\) is the inner automorphism \(x \mapsto axa^{-1}\) of A.

Lemma 2. --- Suppose that B is a semisimple ring that is a finitely generated module over its center Z. Let M and N be (B, A)-bimodules. Suppose that they have finite length (which is, in particular, the case when they are right A-modules of finite length). If M and N are isomorphic as (Z, A)-bimodules, then they are isomorphic as (B, A)-bimodules.

\begin{enumerate}
\def\labelenumi{\Alph{enumi})}
\item
  First consider the case when B is the endomorphism ring of a vector space S of finite dimension d over a commutative field L. We then have Z = L; we view S as a (B, Z)-bimodule. The ring B is simple, S is a simple B-module, and Z is the commutant of S; every B-module is isotypical of type S (VIII, p.~122, Proposition 2 a)). Let \((V, \alpha)\) (resp. \((W, \beta)\) ) be a description of the B-module M (resp. N). The set V (resp. W) is endowed with a (Z, A)-bimodule structure such that \(\alpha\) (resp. \(\beta\) ) is an isomorphism of (B, A)-bimodules (VIII, p.~64, Remark 2). As (Z, A)-bimodules, M is isomorphic to \(V^{d}\) and N to \(W^{d}\) , and there exists an isomorphism from the set of (Z, A)-sub-bimodules of V, ordered by inclusion, to that of the (B, A)-sub-bimodules of M (loc. cit.). Hence V is a (Z, A)-bimodule of finite length, and so is W. Since the (Z, A)-bimodules \(V^{d}\) and \(W^{d}\) are isomorphic, the (Z, A)-bimodules V and W are isomorphic by Theorem 2, d) of VIII, p.~37 applied to the ring \(Z \otimes_{Z} A^{\circ}\) . Finally, the (B, A)-bimodules M and N are isomorphic.
\item
  Now consider the case when B is a simple ring that is finitely generated as a Z-module. Then Z is a field, and B is a central simple algebra of finite degree over the field Z. By Theorem 1 of VIII, p.~252, there exists an extension \(Z'\) of Z of finite degree over Z such that the \(Z'\) -algebra \(B' = B_{(Z')}\) is isomorphic to the endomorphism algebra of a finite-dimensional \(Z'\) -vector space. Set \(M' = M_{(Z')}\) and \(N' = N_{(Z')}\) . Then \(M'\) and \(N'\) are \((B', A)\) -bimodules of finite length; viewed as \((Z', A)\) -bimodules, \(M'\) and \(N'\) are isomorphic. By the case treated in A), \(M'\) and \(N'\) are isomorphic as \((B', A)\) -bimodules and a fortiori as \((B, A)\) -bimodules. Set \(r = [Z' : Z]\) . The \((B, A)\) -bimodule \(M' = Z' \otimes_{Z} M\) is isomorphic to \(M^r\) , and, likewise, the \((B, A)\) -bimodule \(N'\) is isomorphic to \(N^r\) . Since M and N are \((B, A)\) -bimodules of finite length, it follows from Theorem 2, d) of VIII, p.~37 that the \((B, A)\) -bimodules M and N are isomorphic.
\item
  Finally, consider the general case, when B is a semisimple ring that is finitely generated as a Z-module. Let S be the set of classes of simple B-modules; it is finite (VIII, p.~136, Proposition 1). For any \(\lambda \in S\) , denote by \(M_{\lambda}\) (resp. \(N_{\lambda}\) ) the isotypical component of type \(\lambda\) of the B-module M (resp. N); it is a (B, A)-sub-bimodule of M (resp. N) (Remark, VIII, p.~67). For \(\lambda \in S\) , denote the annihilator of the B-module \(\lambda\) by \(b_{\lambda}\) , and set \(B_{\lambda} = B / b_{\lambda}\) ; let \(Z_{\lambda}\) be the center of \(B_{\lambda}\) . For \(\lambda \in S\) , the \((B_{\lambda}, A)\) -bimodules \(M_{\lambda}\) and \(N_{\lambda}\) have finite length. We can then identify B with the product of the simple rings \(B_{\lambda}\) and Z with the product of the \(Z_{\lambda}\) (VIII, p.~141, Proposition 8). Moreover, we can identify M with \(\prod_{\lambda \in S} M_{\lambda}\) and N with \(\prod_{\lambda \in S} N_{\lambda}\) . By assumption, M and N are isomorphic as (Z, A)-bimodules; it follows that for \(\lambda \in S\) , \(M_{\lambda}\) and \(N_{\lambda}\) are isomorphic ( \(Z_{\lambda}, A\) )-bimodules. By the case treated in B), the \((B_{\lambda}, A)\) -bimodules \(M_{\lambda}\) and \(N_{\lambda}\) are isomorphic, and so the (B, A)-bimodules M and N are isomorphic.
\end{enumerate}

Remark. --- It follows from the proof of Lemma 2 that M and N are (Z, A)-bimodules of finite length. Consequently, if B and A are two semisimple rings that are finitely generated modules over their respective centers Z(B) and Z(A), then two (B, A)-bimodules of finite length that are isomorphic as (Z(B), Z(A))-bimodules are isomorphic.
% semantic-fragments: legacy-input-seam
\relax{}
